\chapter{Troubleshooting}
\label{ch:troubleshooting}

When things go wrong, start with the logs and the
\texttt{/admin > Health} dashboard. This chapter catalogues the
problems users actually hit, in rough order of frequency.

\section{Reading the logs}

\begin{lstlisting}
docker compose logs --since 10m mybibli
docker compose logs -f mybibli
\end{lstlisting}

Logs are emitted as structured JSON. The most useful fields are
\texttt{level}, \texttt{target} (which module logged this),
\texttt{message}, and any contextual key=value pairs that the log
line carries (\texttt{user\_id}, \texttt{volume\_id}, \dots).

To filter:

\begin{lstlisting}
docker compose logs --since 10m mybibli | grep -i error
\end{lstlisting}

\section{App does not start}

\subsection*{Symptom: container exits immediately}

Run \texttt{docker compose logs mybibli} and look for one of:

\begin{itemize}
  \item \texttt{missing required environment variable:
        DATABASE\_URL}\index{troubleshooting!DATABASE\_URL} — the
        \texttt{.env} file is missing or the variable is unset. See
        section~\ref{sec:env-vars}.
  \item \texttt{Failed to create database pool} — the database
        hostname / port / credentials in \texttt{DATABASE\_URL} are
        wrong, or the DB is not reachable. From the host, try
        \texttt{docker compose exec mybibli sh -c 'apk add curl;
        curl <db-host>:<db-port>'} to confirm reachability.
  \item \texttt{Failed to run database migrations} — usually a
        permissions issue. Confirm the DB user has \texttt{ALL
        PRIVILEGES} on the database (see
        section~\ref{sec:install-external-db}).
  \item \texttt{Address already in use} — the published host port
        is taken. Change \texttt{HOST\_PORT} in \texttt{.env}.
\end{itemize}

\subsection*{Symptom: every page returns 404}

The container is up but routes are missing. Most common cause:
you pulled an image tag that does not exist (typo) and Docker
fell back to a stale older tag. \texttt{docker compose pull} and
restart.

\section{Cannot log in}

\subsection*{Symptom: ``Invalid credentials'' on a fresh 1.0.x install}

The seed credentials are \texttt{admin/admin}. If they do not work
either:

\begin{itemize}
  \item the database migration that creates them failed silently
        — look for migration errors in the logs;
  \item the previous admin already rotated the password (check
        with whoever installed mybibli);
  \item you are looking at a 1.1.0+ install, where the seed
        credentials are planted by the migration but immediately
        soft-deleted by the issue-\#173 seed gate at startup — use
        the wizard at \texttt{/setup} instead.
\end{itemize}

\section{Stray users in the \texttt{users} table after the wizard}
\label{sec:troubleshoot-seeded-users}

\subsection*{Symptom: four users in the database, only one was created in the wizard}

If you inspect the \texttt{users} table directly (via phpMyAdmin or
the MariaDB CLI) right after completing \texttt{/setup}, you will
see \emph{four} rows even though you only entered one username and
password. This is by design. Three of the rows cannot be used to
log in.

\begin{longtable}{@{} l l c l @{}}
\toprule
\textbf{username} & \textbf{role} & \textbf{active} & \textbf{deleted\_at} \\
\midrule
\endhead
\texttt{admin}     & \texttt{admin}     & 1 & set, $\approx$ first-boot time \\
\texttt{librarian} & \texttt{librarian} & 1 & set, $\approx$ first-boot time \\
\texttt{SYSTEM}    & \texttt{system}    & 0 & NULL \\
\texttt{<your-username>} & \texttt{admin} & 1 & NULL \\
\bottomrule
\end{longtable}

\begin{itemize}
  \item \texttt{admin} and \texttt{librarian} are planted by the
        seed migrations (which run on every fresh install for the
        dev / E2E workflow) and immediately \emph{soft-deleted}
        by the issue-\#173 seed gate at startup. Their
        \texttt{deleted\_at} is set, so the login query (which
        filters \texttt{deleted\_at IS NULL}) refuses them.
  \item \texttt{SYSTEM}\index{SYSTEM user} is a non-loginable
        utility account used to attribute the audit-log entries
        produced by the 30-day auto-purge background task. Its role
        is \texttt{system} (outside the login enum
        \texttt{admin/librarian}), its \texttt{password\_hash} is
        the literal string
        \texttt{NO\_LOGIN\_SYSTEM\_USER\_HASH\_NOT\_A\_VALID\_ARGON2\_STRING},
        and its \texttt{active} flag is \texttt{0}.
\end{itemize}

\begin{tip}
To confirm the soft-deleted rows are inert, try logging in as
\texttt{admin/admin} — it returns an authentication failure. The
\texttt{/admin > Users} panel filters out both the \texttt{SYSTEM}
row and any \texttt{deleted\_at}-set row, so the admin UI shows the
single account you created in the wizard.
\end{tip}

The two soft-deleted rows are physically removed from the table by
the 30-day auto-purge scheduler — expect them to disappear roughly
one month after the boot date stamped in \texttt{deleted\_at}.

\subsection*{Symptom: login redirects back to \texttt{/login}}

The session cookie is being rejected by the browser. Causes:

\begin{itemize}
  \item \texttt{MYBIBLI\_COOKIE\_SECURE=true} on a plain-HTTP
        deployment — the browser refuses to accept a Secure cookie
        over HTTP. Set the variable to \texttt{false} or front the
        instance with HTTPS.
  \item Cookies blocked by browser policy (private window with
        strict third-party blocking, etc.). Test in a default
        profile.
\end{itemize}

\subsection*{Symptom: locked out — last admin deactivated}

The last-active-admin guard should prevent this, but if a previous
admin was deactivated then permanently deleted, recovery requires
direct SQL access:

\begin{lstlisting}
docker compose exec db mariadb -u root -p"$MYSQL_ROOT_PASSWORD" mybibli \
    -e "UPDATE users SET deleted_at=NULL, active=1 WHERE username='your-admin';"
\end{lstlisting}

Then log in normally.

\section{Metadata never resolves}

\subsection*{Symptom: skeleton title stays unresolved}

Open \texttt{/admin > Health}. Look at each provider's status. A
red cross usually means one of:

\begin{itemize}
  \item the provider is genuinely down (rare for BnF / Open Library;
        more common for the smaller services);
  \item the API key is missing or invalid — open
        \texttt{/admin > System > Metadata Providers} and re-enter
        the key;
  \item the host cannot reach the upstream. Confirm with
        \texttt{docker compose exec mybibli sh -c "wget -q
        https://www.googleapis.com/books/v1 -O -"} or equivalent.
\end{itemize}

\subsection*{Symptom: only French ISBNs resolve}

You disabled (or never set) the Google Books / Open Library keys.
BnF is excellent for French-published titles but limited for
foreign-language titles. Add at least Open Library (no key
required).

\section{Search returns too little}

\subsection*{Symptom: known titles do not surface}

\begin{itemize}
  \item Confirm the title is not soft-deleted (check
        \texttt{/admin > Trash}). Search hides soft-deleted items.
  \item Confirm the title's contributors are spelled the way you
        searched. mybibli's search is fuzzy but not magical;
        ``Camus'' will find ``Albert Camus'' but ``Camu'' might
        not.
\end{itemize}

\subsection*{Symptom: search is slow}

For collections above 10,000 titles, queries that combine
contributor + free-text can take a second or two. This is a known
limitation; planned indexes will close the gap.

\section{Cover images are missing}

\subsection*{Symptom: every title shows the placeholder}

The \texttt{covers/} directory is empty or unreadable. Confirm:

\begin{itemize}
  \item the volume mount in \texttt{docker-compose.yml} points to
        a directory that exists and is writable by the container;
  \item the metadata fetch actually returned a cover URL (check the
        title page — if the URL is missing, the provider did not
        supply one).
\end{itemize}

\section{Where to file bugs}

Issues you cannot resolve, or behaviour that contradicts this
manual: open a ticket at
\url{https://github.com/guycorbaz/mybibli/issues}. The issue
template asks for the version (visible in the page footer), the
install method, and the steps to reproduce — please fill them in.
